%! Logarithmic Expressions — Unit 5, Lesson 5.1 — Guided Notes (Answer Key)
\documentclass[10pt]{article}
\usepackage{precalculus-article}
\usepackage{precalculus-key}   % pulls in -boxes; do NOT also load -boxes
\newcommand{\TallMath}[1]{$\displaystyle #1\rule[-1.4em]{0pt}{3.2em}$}

\begin{document}

\pageheader{Unit 5, Lesson 5.1}{Guided Notes --- Answer Key}
\namedateperiod

\vspace{0.04in}

\begin{objectivebox}
By the end of this lesson, I will be able to\ldots
\begin{enumerate}[itemsep=2pt]
  \item Convert between \ans{exponential} form and \ans{logarithmic} form. \hfill\textit{(LO 2.9.A)}
  \item Evaluate logarithmic expressions exactly and estimate others with technology. \hfill\textit{(LO 2.9.A)}
  \item Interpret a \ans{logarithmic} scale, recognizing that each unit is a multiplicative change. \hfill\textit{(LO 2.9.A)}
\end{enumerate}
\end{objectivebox}

\vspace{0.04in}

\begin{vocabbox}
\termblanklong{Logarithm}
\termblanklong{Common logarithm}
\termblanklong{Natural logarithm}
\termblanklong{Exponential form}
\termblanklong{Logarithmic form}
\termblanklong{Logarithmic scale}
\end{vocabbox}

\vspace{0.04in}

\begin{hookbox}
The Richter scale measures earthquake intensity. Magnitude 6 is \textbf{10 times} stronger
than magnitude 5. Magnitude 7 is \ans{100} times stronger than magnitude 5.

Each step on the Richter scale multiplies the intensity by \ans{10}.
This is a \ans{logarithmic} scale --- equal spacing means equal \ans{ratios}.
\end{hookbox}

\vspace{0.04in}

\begin{notesbox}{LO 2.9.A Part 1 --- Definition and Conversion}
\textbf{Definition:}
\[
  \log_b c = a \quad\Leftrightarrow\quad b^a = c
  \qquad b > 0,\ b \ne 1,\ c > 0
\]
``$\log_b c$'' is read: ``log base \ans{$b$} of \ans{$c$}'', and it equals the
\ans{exponent} to which we raise $b$ to get $c$.

\medskip
\textbf{Conversion table} --- fill in the missing logarithmic form:

\renewcommand{\arraystretch}{1.8}
\begin{tabular}{|c|c|}
\hline
\textbf{Exponential Form} & \textbf{Logarithmic Form} \\
\hline
$2^3 = 8$       & $\log_2 8 = $ \ans{3} \\
\hline
$5^2 = 25$      & $\log_5 25 = $ \ans{2} \\
\hline
$10^4 = 10{,}000$ & $\log_{10} 10{,}000 = $ \ans{4} \\
\hline
$b^0 = 1$       & $\log_b 1 = $ \ans{0} \\
\hline
$b^1 = b$       & $\log_b b = $ \ans{1} \\
\hline
\end{tabular}

\medskip
\textbf{Evaluate without a calculator:}

\medskip
$\log_3 81 = $ \ans{4} \qquad $\log_4 2 = $ \ans{$1/2$} \qquad $\log_{10} 0.01 = $ \ans{$-2$}
\end{notesbox}

\vspace{0.04in}

\begin{notesbox}{LO 2.9.A Part 2 --- Common Log and Natural Log}
\textbf{Common logarithm:} $\log x = \log_{\ans{10}} x$ \quad (base 10 when no base is written)

\medskip
\textbf{Natural logarithm:} $\ln x = \log_{\ans{$e$}} x$ \quad where $e \approx $ \ans{2.718}

\medskip
\textbf{Exact values:} $\ln e = $ \ans{1} \qquad $\ln 1 = $ \ans{0}
\qquad $\log_{10} 1 = $ \ans{0}

\medskip
\textbf{Technology estimates:} $\log_2 7 \approx $ \ans{2.807}
\quad (check: $2^{\ans{2}} < 7 < 2^{\ans{3}}$, so the answer is between \ans{2} and \ans{3})
\end{notesbox}

\vspace{0.04in}

\begin{notesbox}{LO 2.9.A Part 3 --- Logarithmic Scale}
On a base-10 logarithmic scale, positions are:
\[
  10^0 = 1 \qquad 10^1 = 10 \qquad 10^2 = 100 \qquad 10^3 = 1{,}000
\]
Each unit represents a \ans{multiplicative} change by \ans{10}.

\medskip
What number sits halfway between $10^1$ and $10^2$ \textit{on the log scale}?

\medskip
Halfway in exponent: $10^{1.5} = $ \ans{$\approx 31.6$} (not 55, the arithmetic midpoint).

\medskip
\textbf{Key idea:} ``Halfway'' on a log scale means the \ans{geometric} mean, not the arithmetic mean.
\end{notesbox}

\end{document}
